\documentclass[11pt]{article}
\usepackage[a4paper,margin=22mm]{geometry}
\usepackage{fontspec}
\setmainfont{DejaVu Serif}
\setsansfont{DejaVu Sans}
\usepackage{microtype}
\usepackage{xcolor}
\usepackage{hyperref}
\usepackage{enumitem}
\usepackage{parskip}
\definecolor{Accent}{HTML}{971D32}
\definecolor{Muted}{HTML}{666666}
\hypersetup{colorlinks=true,urlcolor=Accent,linkcolor=Accent}
\setlist{leftmargin=1.6em,itemsep=0.3em,topsep=0.35em}
\setcounter{tocdepth}{2}
\renewcommand{\contentsname}{Contents}
\newcommand{\sectionrule}{\par\vspace{-0.5em}\noindent\textcolor{Accent}{\rule{\linewidth}{0.6pt}}\par}
\begin{document}

\begin{center}
{\sffamily\bfseries\color{Accent} TOEFL WORD OF THE DAY\par}
\vspace{8mm}
{\fontsize{42}{48}\selectfont\bfseries echelon\par}
\vspace{2mm}
{\Large\sffamily\color{Accent} /ˈɛʃəˌlɑːn/\par}
\vspace{2mm}
{\small\sffamily Local audio phrase: echelon\par}
{\small\sffamily Audio file: audios/echelon.mp3\par}
\vspace{2mm}
{\large\sffamily\color{Muted} countable noun\par}
\vspace{5mm}
{\sffamily a level of authority or status \textbullet\ a diagonal military formation\par}
\vspace{6mm}
{\large An echelon is either a level within a hierarchy or a formation arranged in a series of offset positions.\par}
\vspace{6mm}
\begin{minipage}{0.84\textwidth}
\itshape “The proposal must be approved by the upper echelons of management.”
\end{minipage}\par
\vspace{10mm}
\textbf{Date:} July 28th 2026\quad\textbullet\quad \textbf{Level:} B1--mid-B2 TOEFL
\vfill
Created and compiled by \textbf{Amir Kooshky}\par
\vspace{2mm}
\href{https://t.me/KooshkyTOEFL}{Telegram Channel}\quad\textbullet\quad
\href{https://www.instagram.com/kooshkytoefl}{Instagram}\quad\textbullet\quad
\href{https://t.me/KooshkyTOEFL\_pv}{Personal Telegram}
\end{center}
\newpage
\tableofcontents
\newpage

\section{Meanings and usage}
\sectionrule
\subsection{Meaning 1: Level of authority or status}
\textbf{Part of speech:} countable noun

\textbf{IPA:} /ˈɛʃəˌlɑːn/

\textbf{Pronunciation phrase:} echelon

\textbf{Local audio file:} audios/echelon.mp3

\textbf{Register:} formal; common in business, politics, sociology, and institutional discussions

\textbf{Definition:} a particular level of authority, importance, or status within an organization, profession, or society

\textbf{In simpler words:} a level in a hierarchy

\textbf{Typical pattern:} the upper, higher, or top echelons of + organization or field

\subsubsection{Natural examples}
\begin{enumerate}
  \item The proposal must be approved by the upper echelons of management.
  \item Women remain underrepresented in the highest echelons of the industry.
  \item Information often moves slowly between different echelons of a large organization.
  \item Only officials in the top echelon had access to the confidential report.
  \item The reforms affected employees at every echelon of the institution.
  \item He gradually rose through the echelons of the diplomatic service.
  \item The researcher interviewed people from several social echelons.
\end{enumerate}

\subsubsection{Nuance and implication}
\textbf{Central nuance:} The word presents an organization or society as a hierarchy with several levels.

\textbf{Usual implication:} It is especially common when discussing people with high authority, influence, or status.

\textbf{Compared with rank:} Rank usually refers to the official position of one person. Echelon refers more broadly to a whole level or group within a hierarchy.

\subsubsection{Grammar notes}
\textbf{How to use it:} Use an echelon for one level and echelons for several levels.

\textbf{What can follow it:} It is often followed by of and the name of an organization, profession, industry, or society.

\textbf{Common preposition:} Use in when describing someone within a level, as in people in the highest echelons of government.

\textbf{Noun use:} This is a countable noun. The plural form is especially common in expressions such as the upper echelons and the highest echelons.

\subsubsection{Useful patterns}
\textbf{Pattern:} the upper echelons of + organization or field

\textbf{Use:} Use this for the people at the highest and most powerful levels.

\textbf{Example:} The policy was debated within the upper echelons of the company.

\textbf{Pattern:} the highest or top echelon

\textbf{Use:} Use this for one highest level of authority, ability, or status.

\textbf{Example:} Few athletes ever reach the top echelon of international competition.

\textbf{Pattern:} at every echelon of + organization

\textbf{Use:} Use this when something affects or involves all levels.

\textbf{Example:} The new safety rules apply at every echelon of the organization.

\textbf{Pattern:} rise through the echelons of + profession

\textbf{Use:} Use this when someone gradually moves to increasingly important positions.

\textbf{Example:} She rose through the echelons of the civil service.

\subsubsection{Common wrong usage}
\paragraph{Correction 1}
\textbf{Wrong:} The decision was made by the high echelon of the company.

\textbf{Correct:} The decision was made by the upper echelons of the company.

\textbf{Why:} Upper echelons is the usual expression for the most powerful levels of an organization.

\paragraph{Correction 2}
\textbf{Wrong:} He works on the highest echelon of government.

\textbf{Correct:} He works in the highest echelon of government.

\textbf{Why:} Use in when referring to someone who belongs to a particular level.

\paragraph{Correction 3}
\textbf{Wrong:} The company has several different echelon levels.

\textbf{Correct:} The company has several different echelons.

\textbf{Why:} Echelon already means a level, so echelon level is usually unnecessary.

\subsection{Meaning 2: Diagonal military formation}
\textbf{Part of speech:} countable noun

\textbf{IPA:} /ˈɛʃəˌlɑːn/

\textbf{Pronunciation phrase:} echelon

\textbf{Local audio file:} audios/echelon.mp3

\textbf{Register:} specialized but common in military and historical descriptions

\textbf{Definition:} a military formation in which each group of soldiers, vehicles, ships, or aircraft is positioned behind and to one side of the group ahead

\textbf{In simpler words:} a diagonal line of military units

\textbf{Typical pattern:} in echelon or an echelon formation

\subsubsection{Natural examples}
\begin{enumerate}
  \item The aircraft approached the target in an echelon formation.
  \item The ships moved in echelon to maintain a clear view of one another.
  \item The soldiers advanced in echelon across the open ground.
  \item The vehicles formed a right echelon before entering the valley.
  \item An echelon arrangement allowed each unit to support the one beside it.
  \item The pilots practiced changing from a straight line into an echelon.
  \item The commander ordered the troops to move forward in a left echelon.
\end{enumerate}

\subsubsection{Nuance and implication}
\textbf{Central nuance:} The units are not directly beside or behind one another; each is offset to the left or right.

\textbf{Usual implication:} This arrangement can improve visibility, coordination, or the ability of units to support one another.

\subsubsection{Grammar notes}
\textbf{How to use it:} Use an echelon for the formation itself. The expression in echelon describes how units are arranged or moving.

\textbf{What can follow it:} The noun is commonly used before formation, or after words such as form, move, advance, and fly.

\textbf{Position in a sentence:} In the phrase echelon formation, echelon comes before the noun formation.

\subsubsection{Useful patterns}
\textbf{Pattern:} an echelon formation

\textbf{Use:} Use this to name the diagonal arrangement.

\textbf{Example:} The helicopters flew in an echelon formation.

\textbf{Pattern:} move, advance, or fly in echelon

\textbf{Use:} Use this to describe units traveling in the formation.

\textbf{Example:} The armored vehicles advanced in echelon.

\textbf{Pattern:} a left or right echelon

\textbf{Use:} Use left or right to show the direction in which the formation extends.

\textbf{Example:} The squadron shifted into a right echelon.

\subsubsection{Common wrong usage}
\paragraph{Correction 1}
\textbf{Wrong:} The aircraft flew with echelon.

\textbf{Correct:} The aircraft flew in echelon.

\textbf{Why:} Use in echelon to describe movement in this formation.

\paragraph{Correction 2}
\textbf{Wrong:} The troops formed into an echelon formation.

\textbf{Correct:} The troops formed an echelon.

\textbf{Why:} Form an echelon is more direct. You can also say moved into an echelon formation.

\section{Common collocations}
\sectionrule
\subsection{the upper echelons}
\textbf{Applies to:} Level of authority or status

\textbf{Meaning:} the highest and most powerful levels

\textbf{Pattern:} the upper echelons of + organization or field

\textbf{Example:} The scandal reached the upper echelons of the administration.

\subsection{the highest echelons}
\textbf{Applies to:} Level of authority or status

\textbf{Meaning:} the levels with the greatest authority, influence, or status

\textbf{Pattern:} the highest echelons of + organization, profession, or society

\textbf{Example:} Graduates of the program have entered the highest echelons of public service.

\subsection{top echelon}
\textbf{Applies to:} Level of authority or status

\textbf{Meaning:} the highest level of quality, power, or status

\textbf{Pattern:} the top echelon of + field or group

\textbf{Example:} The university belongs to the top echelon of research institutions.

\subsection{lower echelons}
\textbf{Applies to:} Level of authority or status

\textbf{Meaning:} levels with less authority or status

\textbf{Pattern:} the lower echelons of + organization

\textbf{Example:} The policy was unpopular among workers in the lower echelons of the organization.

\subsection{every echelon of}
\textbf{Applies to:} Level of authority or status

\textbf{Meaning:} every level within a hierarchy

\textbf{Pattern:} at every echelon of + organization

\textbf{Example:} The investigation revealed problems at every echelon of the institution.

\subsection{echelon formation}
\textbf{Applies to:} Diagonal military formation

\textbf{Meaning:} a formation in which each unit is behind and to one side of another

\textbf{Example:} The aircraft maintained an echelon formation during the exercise.

\subsection{move in echelon}
\textbf{Applies to:} Diagonal military formation

\textbf{Meaning:} travel while arranged in a diagonal formation

\textbf{Pattern:} move, advance, or fly in echelon

\textbf{Example:} The ships moved in echelon along the coast.

\subsection{right echelon}
\textbf{Applies to:} Diagonal military formation

\textbf{Meaning:} a formation extending diagonally toward the right

\textbf{Example:} The vehicles formed a right echelon as they approached the bridge.

\section{Synonyms and antonyms}
\sectionrule
\subsection{Close synonyms}
\subsubsection{level}
\textbf{Part of speech:} countable noun

\textbf{Applies to:} Level of authority or status

\textbf{Shared meaning:} Both words can refer to a position within a hierarchy.

\textbf{Difference:} Level is a broad everyday word. Echelon is more formal and emphasizes a position within an organized hierarchy.

\textbf{Example:} The issue must be discussed at the highest level of government.

\subsubsection{tier}
\textbf{Part of speech:} countable noun

\textbf{Applies to:} Level of authority or status

\textbf{Shared meaning:} Both words describe one level in a structure containing several levels.

\textbf{Difference:} Tier commonly describes one of several clearly separated levels in a system, service, or structure. Echelon is especially associated with authority, status, or organizational hierarchy.

\textbf{Example:} Customers can choose from three different service tiers.

\subsubsection{rank}
\textbf{Part of speech:} countable noun

\textbf{Applies to:} Level of authority or status

\textbf{Shared meaning:} Both can indicate a position within a hierarchy.

\textbf{Difference:} Rank often identifies the official position of a particular person. Echelon usually describes a broader level containing several people.

\textbf{Example:} She achieved the rank of senior officer.

\subsubsection{formation}
\textbf{Part of speech:} countable noun

\textbf{Applies to:} Diagonal military formation

\textbf{Shared meaning:} Both words can describe an organized arrangement of soldiers, vehicles, ships, or aircraft.

\textbf{Difference:} Formation is a general word for any organized arrangement of military units. Echelon refers to one specific diagonal arrangement.

\textbf{Example:} The pilots maintained a tight formation.

\section{Etymology}
\sectionrule
\textbf{Origin:} The word entered English from French échelon, meaning a rung of a ladder or a step.

\section{Practice exercises}
\sectionrule
\subsection{Correct or incorrect?}
Decide whether each sentence is correct. Correct the inaccurate ones.

\begin{enumerate}
  \item The policy was approved by officials in the upper echelons of government.
  \item She works on the highest echelon of the organization.
  \item The company has four separate echelon levels.
  \item The aircraft flew in echelon during the military exercise.
\end{enumerate}

\subsection{Collocation practice}
Complete each sentence with a natural collocation.

\begin{enumerate}
  \item The investigation involved officials in the upper \_\_\_\_\_ of the ministry.
  \item Very few musicians reach the top \_\_\_\_\_ of the profession.
  \item The helicopters flew in an echelon \_\_\_\_\_ during the demonstration.
\end{enumerate}

\subsection{Complete answer key}
\subsubsection{Correct or incorrect?}
\begin{enumerate}
  \item \textbf{Correct} \emph{The upper echelons of is a natural expression for the highest levels of authority.}
  \item \textbf{Incorrect}. She works in the highest echelon of the organization. \emph{Use in when someone belongs to a particular echelon.}
  \item \textbf{Incorrect}. The company has four separate echelons. \emph{Echelon already means level, so adding levels is redundant.}
  \item \textbf{Correct} \emph{In echelon correctly describes movement in a diagonal formation.}
\end{enumerate}

\subsubsection{Collocation practice}
\begin{enumerate}
  \item \textbf{echelons.} The investigation involved officials in the upper echelons of the ministry. \emph{The upper echelons is the standard expression for the most powerful levels.}
  \item \textbf{echelon.} Very few musicians reach the top echelon of the profession. \emph{Top echelon means the highest level of achievement or status.}
  \item \textbf{formation.} The helicopters flew in an echelon formation during the demonstration. \emph{Echelon formation is the established term for this diagonal arrangement.}
\end{enumerate}

\vfill
\begin{center}
\small\color{Muted} Created and compiled by \textbf{Amir Kooshky}\\
\href{https://t.me/KooshkyTOEFL}{Telegram Channel}\quad\textbullet\quad
\href{https://www.instagram.com/kooshkytoefl}{Instagram}\quad\textbullet\quad
\href{https://t.me/KooshkyTOEFL\_pv}{Personal Telegram}
\end{center}
\end{document}
